% Options for packages loaded elsewhere
\PassOptionsToPackage{unicode}{hyperref}
\PassOptionsToPackage{hyphens}{url}
\PassOptionsToPackage{dvipsnames,svgnames,x11names}{xcolor}
%
\documentclass[
  letterpaper,
  DIV=11,
  numbers=noendperiod]{scrreprt}

\usepackage{amsmath,amssymb}
\usepackage{iftex}
\ifPDFTeX
  \usepackage[T1]{fontenc}
  \usepackage[utf8]{inputenc}
  \usepackage{textcomp} % provide euro and other symbols
\else % if luatex or xetex
  \usepackage{unicode-math}
  \defaultfontfeatures{Scale=MatchLowercase}
  \defaultfontfeatures[\rmfamily]{Ligatures=TeX,Scale=1}
\fi
\usepackage{lmodern}
\ifPDFTeX\else  
    % xetex/luatex font selection
\fi
% Use upquote if available, for straight quotes in verbatim environments
\IfFileExists{upquote.sty}{\usepackage{upquote}}{}
\IfFileExists{microtype.sty}{% use microtype if available
  \usepackage[]{microtype}
  \UseMicrotypeSet[protrusion]{basicmath} % disable protrusion for tt fonts
}{}
\makeatletter
\@ifundefined{KOMAClassName}{% if non-KOMA class
  \IfFileExists{parskip.sty}{%
    \usepackage{parskip}
  }{% else
    \setlength{\parindent}{0pt}
    \setlength{\parskip}{6pt plus 2pt minus 1pt}}
}{% if KOMA class
  \KOMAoptions{parskip=half}}
\makeatother
\usepackage{xcolor}
\setlength{\emergencystretch}{3em} % prevent overfull lines
\setcounter{secnumdepth}{-\maxdimen} % remove section numbering
% Make \paragraph and \subparagraph free-standing
\makeatletter
\ifx\paragraph\undefined\else
  \let\oldparagraph\paragraph
  \renewcommand{\paragraph}{
    \@ifstar
      \xxxParagraphStar
      \xxxParagraphNoStar
  }
  \newcommand{\xxxParagraphStar}[1]{\oldparagraph*{#1}\mbox{}}
  \newcommand{\xxxParagraphNoStar}[1]{\oldparagraph{#1}\mbox{}}
\fi
\ifx\subparagraph\undefined\else
  \let\oldsubparagraph\subparagraph
  \renewcommand{\subparagraph}{
    \@ifstar
      \xxxSubParagraphStar
      \xxxSubParagraphNoStar
  }
  \newcommand{\xxxSubParagraphStar}[1]{\oldsubparagraph*{#1}\mbox{}}
  \newcommand{\xxxSubParagraphNoStar}[1]{\oldsubparagraph{#1}\mbox{}}
\fi
\makeatother


\providecommand{\tightlist}{%
  \setlength{\itemsep}{0pt}\setlength{\parskip}{0pt}}\usepackage{longtable,booktabs,array}
\usepackage{calc} % for calculating minipage widths
% Correct order of tables after \paragraph or \subparagraph
\usepackage{etoolbox}
\makeatletter
\patchcmd\longtable{\par}{\if@noskipsec\mbox{}\fi\par}{}{}
\makeatother
% Allow footnotes in longtable head/foot
\IfFileExists{footnotehyper.sty}{\usepackage{footnotehyper}}{\usepackage{footnote}}
\makesavenoteenv{longtable}
\usepackage{graphicx}
\makeatletter
\def\maxwidth{\ifdim\Gin@nat@width>\linewidth\linewidth\else\Gin@nat@width\fi}
\def\maxheight{\ifdim\Gin@nat@height>\textheight\textheight\else\Gin@nat@height\fi}
\makeatother
% Scale images if necessary, so that they will not overflow the page
% margins by default, and it is still possible to overwrite the defaults
% using explicit options in \includegraphics[width, height, ...]{}
\setkeys{Gin}{width=\maxwidth,height=\maxheight,keepaspectratio}
% Set default figure placement to htbp
\makeatletter
\def\fps@figure{htbp}
\makeatother

\KOMAoption{captions}{tableheading}
\makeatletter
\@ifpackageloaded{caption}{}{\usepackage{caption}}
\AtBeginDocument{%
\ifdefined\contentsname
  \renewcommand*\contentsname{Table of contents}
\else
  \newcommand\contentsname{Table of contents}
\fi
\ifdefined\listfigurename
  \renewcommand*\listfigurename{List of Figures}
\else
  \newcommand\listfigurename{List of Figures}
\fi
\ifdefined\listtablename
  \renewcommand*\listtablename{List of Tables}
\else
  \newcommand\listtablename{List of Tables}
\fi
\ifdefined\figurename
  \renewcommand*\figurename{Figure}
\else
  \newcommand\figurename{Figure}
\fi
\ifdefined\tablename
  \renewcommand*\tablename{Table}
\else
  \newcommand\tablename{Table}
\fi
}
\@ifpackageloaded{float}{}{\usepackage{float}}
\floatstyle{ruled}
\@ifundefined{c@chapter}{\newfloat{codelisting}{h}{lop}}{\newfloat{codelisting}{h}{lop}[chapter]}
\floatname{codelisting}{Listing}
\newcommand*\listoflistings{\listof{codelisting}{List of Listings}}
\makeatother
\makeatletter
\makeatother
\makeatletter
\@ifpackageloaded{caption}{}{\usepackage{caption}}
\@ifpackageloaded{subcaption}{}{\usepackage{subcaption}}
\makeatother

\ifLuaTeX
  \usepackage{selnolig}  % disable illegal ligatures
\fi
\usepackage{bookmark}

\IfFileExists{xurl.sty}{\usepackage{xurl}}{} % add URL line breaks if available
\urlstyle{same} % disable monospaced font for URLs
\hypersetup{
  colorlinks=true,
  linkcolor={blue},
  filecolor={Maroon},
  citecolor={Blue},
  urlcolor={Blue},
  pdfcreator={LaTeX via pandoc}}


\author{}
\date{}

\begin{document}


\chapter{Documentation Outline}\label{documentation-outline}

Moved to Hedgedoc 28 May 2024 (\textbf{mrchristian?}) :
https://demo.hedgedoc.org/QyniBhytQGGeIx3XdM8F2g

The idea is to create descriptions of what we would like the ALR to do.

The writing style will be based on the following:

The guide will be considered a Step-by-step guide acording to the
Diátaxis method as the audience will be users familiar with data science
tooling and paper writing (as Diátaxis describes it - a work place
instruction).

\begin{itemize}
\tightlist
\item
  Technical writing method: Diátaxis - https://diataxis.fr/
\item
  Technical writing style: MDN Web Docs -
  https://developer.mozilla.org/en-US/docs/MDN/Writing\_guidelines/Writing\_style\_guide
\item
  Style guide: APA - https://apastyle.apa.org/
\end{itemize}

The writing tooling has not yet been selected.

Contributors:

\begin{itemize}
\tightlist
\item
  \href{https://github.com/neerajkumaris}{(\textbf{neerajkumaris?})}
\item
  \href{https://github.com/ydv-amit-ydv}{(\textbf{Amit?})}
\item
  \href{https://github.com/Adesh212}{(\textbf{Adesh212?})}
\item
  \href{https://github.com/mrchristian}{(\textbf{Simon?}) Worthington}
\end{itemize}

\section{SECTIONS}\label{sections}

\section{Orientation and learning
points}\label{orientation-and-learning-points}

Assigned (\textbf{NOT?}) ASSIGNED -

\textbf{Questions:}

\textbf{Notes:}

\textbf{Outline:}

\section{A literature review
question?}\label{a-literature-review-question}

Assigned
\href{https://github.com/neerajkumaris}{(\textbf{neerajkumaris?})} -

\textbf{Questions:}

How does the author report on their question in the literature review in
terms of the use of ALR?

\textbf{Notes:}

When the author reports on their question in the literature review using
Automated Literature Review (ALR) tools such as pygetpapers, they
typically follow a systematic process that leverages the capabilities of
the tool to enhance the efficiency and comprehensiveness of their
review.

\textbf{Outline:} Outline for Reporting on the Research Question in the
Literature Review Using pygetpapers :

\textbf{Introduction:}

\begin{enumerate}
\def\labelenumi{\arabic{enumi}.}
\tightlist
\item
  Introduce the research question and its significance.
\item
  Provide background information on the topic.
\item
  Explain the relevance of Automated Literature Review (ALR) tools like
  pygetpapers.
\item
  State the purpose of the literature review and its organization.
\end{enumerate}

\textbf{Methodology for Literature Search:}

\begin{enumerate}
\def\labelenumi{\arabic{enumi}.}
\tightlist
\item
  Describe the methodology used to conduct the literature search.
\item
  Explain the use of pygetpapers for automating the retrieval of
  academic papers.
\item
  Specify the search strategy, including keywords, phrases, and Boolean
  operators.
\item
  Detail the sources from which papers were retrieved (e.g., Europe
  PMC).
\end{enumerate}

\textbf{Selection Criteria and Filtering:}

\begin{enumerate}
\def\labelenumi{\arabic{enumi}.}
\tightlist
\item
  Define the inclusion and exclusion criteria used to filter the search
  results.
\item
  Discuss the process of filtering papers based on relevance and
  quality.
\item
  Highlight any manual or programmatic filtering techniques employed.
\end{enumerate}

\textbf{Organization and Categorization:}

\begin{enumerate}
\def\labelenumi{\arabic{enumi}.}
\tightlist
\item
  Explain how the retrieved literature was organized and categorized.
\item
  Discuss the criteria used for categorization, such as themes,
  methodologies, or findings.
\item
  Provide insights into the structure of the literature review based on
  these categories.
\end{enumerate}

\textbf{Synthesis of Findings:}

\begin{enumerate}
\def\labelenumi{\arabic{enumi}.}
\tightlist
\item
  Summarize the main findings from the retrieved literature.
\item
  Discuss common themes, trends, and patterns identified in the
  literature.
\item
  Highlight any contradictions or inconsistencies among the findings.
\end{enumerate}

\textbf{Critical Analysis and Quality Assessment:}

\begin{enumerate}
\def\labelenumi{\arabic{enumi}.}
\tightlist
\item
  Conduct a critical analysis of the methodologies employed in the
  retrieved studies.
\item
  Assess the reliability, validity, and generalizability of the
  findings.
\item
  Utilize metrics such as citation counts or impact factors to evaluate
  the quality of the literature.
\end{enumerate}

\textbf{Conclusion:}

\begin{enumerate}
\def\labelenumi{\arabic{enumi}.}
\tightlist
\item
  Summarize the key findings and insights from the literature review.
\item
  Reiterate the significance of the research question and its
  implications for the field.
\item
  Reflect on the role of ALR tools like pygetpapers in facilitating the
  literature review process.
\end{enumerate}

\section{Using CoLab}\label{using-colab}

Assigned
\href{https://github.com/neerajkumaris}{(\textbf{neerajkumaris?})} -

\textbf{Questions:}

\begin{itemize}
\tightlist
\item
  How and when to add CoLab to GitHub?
  \href{https://github.com/mrchristian}{(\textbf{Simon?}) Worthington}
\end{itemize}

\textbf{Notes:}

\textbf{Outline:}

\textbf{When to Add CoLab to GitHub}

\begin{enumerate}
\def\labelenumi{\arabic{enumi}.}
\item
  To keep track of changes and revisions to your notebooks.
\item
  To easily share notebooks with collaborators and work together on the
  same project.
\item
  To keep a backup of your work in case of accidental deletions or
  changes.
\item
  To integrate your notebooks into a larger project managed on GitHub,
  benefiting from issues tracking, pull requests, and other GitHub
  features.
\end{enumerate}

\textbf{Steps to integrate Google Colab to GitHub}

\textbf{1. Linking GitHub to Colab}

\textbf{First, you need to link your GitHub account with your Google
Colab account:}

\begin{enumerate}
\def\labelenumi{\alph{enumi}.}
\item
  Open a Google Colab notebook.
\item
  Click on ``File'' in the menu.
\item
  Select ``Save a copy in GitHub''.
\item
  If you haven't linked your GitHub account before, a prompt will appear
  asking you to authenticate with GitHub. Follow the instructions to
  grant access.
\end{enumerate}

\textbf{2. Saving a Colab Notebook to GitHub}

\textbf{Once your GitHub account is linked, you can save notebooks
directly to a GitHub repository:}

\begin{enumerate}
\def\labelenumi{\alph{enumi}.}
\item
  Open the notebook you want to save.
\item
  Click on ``File'' in the menu.
\item
  Select ``Save a copy in GitHub''.
\item
  A dialog will appear where you can specify the repository, the branch,
  and the commit message.
\item
  Click ``OK'' to save the notebook to the specified GitHub repository.
\end{enumerate}

\textbf{3. Opening a Notebook from GitHub in Colab}

\textbf{To open a notebook directly from a GitHub repository:}

\begin{enumerate}
\def\labelenumi{\alph{enumi}.}
\item
  Go to the Google Colab homepage.
\item
  Click on ``GitHub'' in the menu.
\item
  Authenticate with your GitHub account if prompted.
\item
  Search for the repository or notebook you want to open.
\item
  Click on the notebook to open it in Colab.
\end{enumerate}

\textbf{4. Using Git Commands in Colab}

\textbf{For more advanced use cases, you can use Git commands directly
within a Colab notebook:}

\textbf{Start by cloning a repository:}

!git clone https://github.com/yourusername/yourrepository.git

\textbf{Navigate into the repository directory:}

\%cd yourrepository

You can now use other Git commands, such as git pull, git add, git
commit, and git push, just like you would in a local development
environment.

\section{Install}\label{install}

Assigned
\href{https://github.com/neerajkumaris}{(\textbf{neerajkumaris?})} -

\textbf{Questions:}

\textbf{Notes:}

\textbf{Outline:}

\section{Searching open corpora}\label{searching-open-corpora}

Assigned \href{https://github.com/ydv-amit-ydv}{(\textbf{Amit?})}

\subsection{\texorpdfstring{\textbf{Identification of the Topics and
queries}
}{Identification of the Topics and queries  }}\label{identification-of-the-topics-and-queries}

Start by removing the stop words and select the most relevant keywords
For example : Ques : What is the role of silicon transporter in the rice
? \textgreater Step 1 : Remove Stop words : Remaing words are : Role ,
Silicon , transorter, rice

\subsection{\texorpdfstring{\textbf{Searching the corpora (EUPMC)}
}{Searching the corpora (EUPMC)  }}\label{searching-the-corpora-eupmc}

Open the command terminal in ther directory where you wish to store your
project \textgreater Step 2 : \emph{pygetpapers -q ``your query terms''
-n ``lantana\_query\_config''--save\_query}

Using this you will be able to count the total number of available
papers on EUPMC for the query We can also use the attributes like
-startdate and enddate to filter out the query

\begin{quote}
Step 3 : \emph{pygetpapers -q '' role silicon transporter rice '' -k 10
-p -x -makecsv -makehtml -o path/to/some/output/directory/optional
--loglevel debug -x --logfile test\_log.txt}
\end{quote}

This will download the 10 papers from the EUPMC in both the xml and pdf
formats , make the csv and html of the associated metadata and download
these papers in the given output directory Similarly we can download
more papers and make the entire cProject tree
\textgreater{}\emph{--loglevel debug -x --logfile test\_log.txt} is to
write the log to a .txt file in your HOME directory, while
simultaneously printing it out.

\begin{quote}
\emph{When -o output path is not given the downloaded copora is saved to
the pwd as a time-stamed directory}
\end{quote}

\subsection{\texorpdfstring{\textbf{Short Summary of the essential tags
of the
pegetpapers}}{Short Summary of the essential tags of the pegetpapers}}\label{short-summary-of-the-essential-tags-of-the-pegetpapers}

\begin{longtable}[]{@{}
  >{\raggedright\arraybackslash}p{(\columnwidth - 4\tabcolsep) * \real{0.3333}}
  >{\raggedright\arraybackslash}p{(\columnwidth - 4\tabcolsep) * \real{0.3333}}
  >{\raggedright\arraybackslash}p{(\columnwidth - 4\tabcolsep) * \real{0.3333}}@{}}
\toprule\noalign{}
\begin{minipage}[b]{\linewidth}\raggedright
Flag
\end{minipage} & \begin{minipage}[b]{\linewidth}\raggedright
What it does
\end{minipage} & \begin{minipage}[b]{\linewidth}\raggedright
In this case \texttt{pygetpapers}\ldots{}
\end{minipage} \\
\midrule\noalign{}
\endhead
\bottomrule\noalign{}
\endlastfoot
\texttt{-q} & specifies the query & queries for `invasive plant species'
in METHODS section \\
\texttt{-k} & number of hits (default 100) & limits hits to 10 \\
\texttt{-o} & specifies output directory & outputs to
invasive\_plant\_species\_test \\
\texttt{-x} & downloads fulltext xml & \\
\texttt{-c} & saves per-paper metadata into a single csv & saves single
CSV named
\href{resources/invasive_plant_species_test/europe_pmc.csv}{\texttt{europe\_pmc.csv}} \\
\texttt{-\/-makehtml} & saves per-paper metadata into a single HTML file
& saves single HTML named
\href{resources/invasive_plant_species_test/eupmc_results.html}{\texttt{europe\_pmc.html}} \\
\texttt{-\/-save\_query} & saves the given query in a
\texttt{config.ini} in output directory & saves query to
\href{resources/invasive_plant_species_test/saved_config.ini}{\texttt{saved\_config.ini}} \\
\end{longtable}

\texttt{pygetpapers}, by default, writes metadata to a JSON file within:
- individual paper directory for corresponding paper
(\texttt{epmc\_result.json}) - working directory for all downloaded
papers
(\href{resources/invasiv_plant_species_test/eupmc_results.json}{\texttt{epmc\_results.json}})

\subsection{\texorpdfstring{\textbf{Notes :}}{Notes :}}\label{notes}

-The query terms can be also specified in the purticular section like
Abstract, methods, results etc
\textgreater{}C:\Users\lilia\textgreater pygetpapers -q
abstract:``money'' -n INFO: Total number of hits for the query are 22161

\section{Analysing your collection}\label{analysing-your-collection}

Assigned \href{https://github.com/Adesh212}{(\textbf{Adesh212?})} -

\textbf{Questions:}

\textbf{Notes:}

\textbf{Outline:}

\section{Using the literature collection in your review
paper}\label{using-the-literature-collection-in-your-review-paper}

Assigned \href{https://github.com/mrchristian}{(\textbf{Simon?})
Worthington} -

\textbf{Notes:}

The author needs to use the collection, and information and data about
the collection in their literature review paper.

\textbf{Questions:}

What does the collection (or what would we like it to collect) and the
process produce that the author can use in their paper?

\begin{itemize}
\tightlist
\item
  A complete Git repository
\item
  A replicable and reusable system
\item
  A DOI references data and code set on Zenodo or other academic
  repository as a full GitHub repository using Software Citation
\item
  A saved query:
  https://pygetpapers.readthedocs.io/en/latest/user\_documentation.html\#example-query
  Flag --save\_query as saved\_config.ini
\item
  A CProject for the whole query
\item
  A CTree per paper
\item
  Custom terms and dictionaries
\item
  Data analysis results?
\item
  Hits
\item
  Word frequency
\item
  Bibtex output and import with content an/or contnet links to Zotero
\item
  CSS Paged media outputs
\end{itemize}

\textbf{Outline:}

\textbf{Ideal process}

\begin{itemize}
\tightlist
\item
  Add reference list to the paper and store it in Zotero
\item
  Take a GitHub Release of repository add to Zenodo and link as data to
  paper
\item
  Link to all papers in repo CTree
\item
  Present question and methods for converting into query, show table of
  queries, link to data
\item
  Provide information on how to reproduce and reuse / extend the
  CProject
\item
  Present data analysis of literature review in paper
\item
  Chart and data link to number of hits
\item
  Chart of NGram of search term
\item
  Table of word frequency
\item
  Table of top ten papers
\item
  Table of top journals, top authors
\end{itemize}

\textbf{Current process}

\begin{itemize}
\tightlist
\item
  Take a GitHub Release of repository add to Zenodo and link as data to
  paper
\item
  Link to all papers in repo CTree
\item
  Present question and methods for converting into query, show table of
  queries, link to data
\item
  Provide information on how to reproduce and reuse / extend the
  CProject
\item
  Present data analysis of literature review in paper
\item
  Table of the number of hits
\item
  Table of word frequency
\item
  Table of top ten papers
\item
  Table of top journals, top authors
\end{itemize}

\section{Other sections}\label{other-sections}

\subsection{Metadata}\label{metadata}

\subsection{Other publication sections
needed}\label{other-publication-sections-needed}

Covers, title pages, table of contents, contributors, etc. See style
guide: APA - https://apastyle.apa.org/

\subsection{Colophon}\label{colophon}

\subsection{Introduction}\label{introduction}

\subsection{Resources}\label{resources}

\subsection{Glossary and
abbreviations}\label{glossary-and-abbreviations}

\subsection{References}\label{references}

Zotero collection:




\end{document}
